% Source: Weinberg GR, physical PDF pp. 289--295 (printed pp. 267--273).
% Coverage: Section 10.5, Quadrupole Radiation; equations
%           (10.5.1)--(10.5.25).
% Figures/tables/footnotes: reference markers 1--4; no figures or tables.
% Status: source-reviewed and compile-clean.
% Uncertainties: Printed p. 272 says Eq. (10.5.23) gives the point-mass
% radiated power, but (10.5.23) merely defines I; the intended reference is
% evidently Eq. (10.5.22), used below and recorded as a source anomaly.

\subsection{Quadrupole Radiation}\label{sec:10.5}

Up to this point we have made no approximations beyond the basic assumption
that the fields are weak. (Our use of the wave-zone limits \(r\gg R\),
\(r\gg1/\omega\), \(r\gg\omega R^2\) was not really an approximation,
since we can always choose \(r\) large enough to make these assumptions true;
and by the conservation of energy, the power passing through a sphere at
large \(r\) must equal that passing through any surface enclosing the
radiating system.) We now make a further approximation, and assume that the
source radius \(R\) is much smaller than the wavelength \(1/\omega\):
\begin{equation}
  \omega R\ll1.
  \tag{10.5.1}\label{eq:10.5.1}
\end{equation}
Most of the radiation is emitted at frequencies of order
\(\bar v/R\), where \(\bar v\) is some typical velocity within the system,
so we are really making the same sort of approximation as that made in the
previous chapter, that is, \(\bar v\ll1\).

When \eqref{eq:10.5.1} holds we may approximate the Fourier transforms
needed in \eqref{eq:10.4.14} and \eqref{eq:10.4.16} by the
\(\mathbf{k}\)-independent integral
\begin{equation}
  T_{ij}(\mathbf{k},\omega)
  \simeq
  \int \dd^3x\,T_{ij}(\mathbf{x},\omega).
  \tag{10.5.2}\label{eq:10.5.2}
\end{equation}
This can be rewritten in a useful way by using the conservation laws in the
form
\[
  \partial_i\partial_jT^{ij}(\mathbf{x},\omega)
  =
  -\omega^2T^{00}(\mathbf{x},\omega).
\]
Multiplying by \(x^ix^j\) and integrating over \(\mathbf{x}\), we find
\begin{equation}
  T_{ij}(\mathbf{k},\omega)
  \simeq
  -\frac{\omega^2}{2}D_{ij}(\omega),
  \tag{10.5.3}\label{eq:10.5.3}
\end{equation}
\begin{equation}
  D_{ij}(\omega)
  \equiv
  \int \dd^3x\,x^ix^jT^{00}(\mathbf{x},\omega).
  \tag{10.5.4}\label{eq:10.5.4}
\end{equation}
The power per solid angle is therefore
\begin{equation}
  \frac{\dd P}{\dd\Omega}
  =
  \frac{G\omega^6}{4\pi}
  \Lambda_{ij,lm}(\widehat{\mathbf{k}})
  D_{ij}^*(\omega)D_{lm}(\omega).
  \tag{10.5.5}\label{eq:10.5.5}
\end{equation}
If the source is a sum of Fourier components, then the power radiated is a
sum of terms such as \eqref{eq:10.5.5}. If the source is a Fourier integral
like \eqref{eq:10.4.1}, then the energy emitted per solid angle is
\begin{equation}
  \frac{\dd E}{\dd\Omega}
  =
  \frac12G\Lambda_{ij,lm}(\widehat{\mathbf{k}})
  \int_0^\infty
  \omega^6D_{ij}^*(\omega)D_{lm}(\omega)\,\dd\omega.
  \tag{10.5.6}\label{eq:10.5.6}
\end{equation}

The coefficients \(D_{ij}(\omega)\) in \eqref{eq:10.5.5} and
\eqref{eq:10.5.6} do not depend on the direction
\(\widehat{\mathbf{k}}\) of the emitted radiation, so we can do the
integral over solid angle once and for all. We use the formulas
\[
  \int \dd\Omega\,\widehat k_i\widehat k_j
  =
  \frac{4\pi}{3}\delta_{ij},
\]
\[
  \int \dd\Omega\,
  \widehat k_i\widehat k_j\widehat k_l\widehat k_m
  =
  \frac{4\pi}{15}
  \left(
    \delta_{ij}\delta_{lm}
    +\delta_{il}\delta_{jm}
    +\delta_{im}\delta_{jl}
  \right).
\]
(The form of the right-hand sides is dictated by symmetry and rotational
invariance; the numerical coefficients can be calculated by contracting
\(i\) with \(j\) and \(l\) with \(m\).) We then find
\[
  \int \dd\Omega\,
  \Lambda_{ij,lm}(\widehat{\mathbf{k}})
  =
  \frac{2\pi}{15}
  \left(
    11\delta_{il}\delta_{jm}
    -4\delta_{ij}\delta_{lm}
    +\delta_{im}\delta_{jl}
  \right),
\]
so the power emitted at a single discrete frequency \(\omega\) is
\begin{equation}
  P
  =
  \frac{2G\omega^6}{5}
  \left[
    D_{ij}^*(\omega)D_{ij}(\omega)
    -\frac13\left|D_{ii}(\omega)\right|^2
  \right],
  \tag{10.5.7}\label{eq:10.5.7}
\end{equation}
whereas for a smooth distribution of frequencies the total emitted energy is
\begin{equation}
  E
  =
  \frac{4\pi G}{5}
  \int_0^\infty
  \omega^6
  \left[
    D_{ij}^*(\omega)D_{ij}(\omega)
    -\frac13\left|D_{ii}(\omega)\right|^2
  \right]\dd\omega.
  \tag{10.5.8}\label{eq:10.5.8}
\end{equation}

Before going on to calculate the quadrupole radiation emitted in a few
special cases, it will be necessary to pause for a few comments on the method
of calculation:

\medskip
\noindent\textit{(A)} The quadrupole approximation is usually applicable to
nonrelativistic systems, and for these systems the energy density
\(T^{00}(\mathbf{x},\omega)\) is dominated by the rest-mass density of the
system. It may be surprising that we do not need to take explicit account of
the potential and kinetic energy terms in the full tensor \(T^{\mu\nu}\),
because such terms must be included if \(T^{\mu\nu}\) is to be conserved!
Indeed, for a system of particles bound by gravitational forces we should in
principle take \(T^{\mu\nu}\) as the total ``tensor'' \(\tau^{\mu\nu}\)
constructed in Section~\ref{sec:7.6}, including terms nonlinear in the
gravitational fields. However, we have already exploited energy and momentum
conservation in our derivation of Eqs.~\eqref{eq:10.5.3}--\eqref{eq:10.5.6},
and since this has given a result involving only \(T^{00}\), we are now free
to approximate \(T^{00}\) with the rest-mass density.

\medskip
\noindent\textit{(B)} For general systems of vibrating and/or rotating
solids, it is often quite difficult to evaluate the Fourier transform
\(T^{00}(\mathbf{x},\omega)\), defined by Eq.~\eqref{eq:10.4.1} or
\eqref{eq:10.4.2}. It is much easier first to evaluate the moments
\begin{equation}
  D_{ij}(t)
  \equiv
  \int \dd^3x\,x^ix^jT^{00}(\mathbf{x},t),
  \tag{10.5.9}\label{eq:10.5.9}
\end{equation}
and then evaluate \(D_{ij}(\omega)\) by expressing \(D_{ij}(t)\) as a
Fourier integral,
\begin{equation}
  D_{ij}(t)
  =
  \int_0^\infty \dd\omega\,D_{ij}(\omega)e^{-\ii\omega t}
  +\text{c.c.},
  \tag{10.5.10}\label{eq:10.5.10}
\end{equation}
or as a sum of Fourier components,
\begin{equation}
  D_{ij}(t)
  =
  \sum_\omega e^{-\ii\omega t}D_{ij}(\omega)+\text{c.c.}
  \tag{10.5.11}\label{eq:10.5.11}
\end{equation}

\medskip
\noindent\textit{(C)} The question may arise: What origin should be taken for
the coordinates \(x^i\) in the integral \eqref{eq:10.5.4} for \(D_{ij}\)?
In principle, it does not matter. When we shift the origin of coordinates by
an amount \(a^i\), we change \(D_{ij}\) into
\begin{align*}
  \int &(x^i-a^i)(x^j-a^j)T^{00}(\mathbf{x},t)\,\dd^3x\\
  ={}&
  \int x^ix^jT^{00}(\mathbf{x},t)\,\dd^3x
  -a^i\int x^jT^{00}(\mathbf{x},t)\,\dd^3x
  -a^j\int x^iT^{00}(\mathbf{x},t)\,\dd^3x
  +a^ia^j\int T^{00}(\mathbf{x},t)\,\dd^3x.
\end{align*}
But conservation of energy and momentum tells us that the last three terms
are at most linear functions of time, because
\[
  \frac{\partial}{\partial t}\int T^{00}(\mathbf{x},t)\,\dd^3x
  =
  -\int\partial_iT^{i0}(\mathbf{x},t)\,\dd^3x
  =0,
\]
\[
  \frac{\partial^2}{\partial t^2}
  \int x^iT^{00}(\mathbf{x},t)\,\dd^3x
  =
  \int x^i\partial_j\partial_kT^{jk}(\mathbf{x},t)\,\dd^3x
  =
  -\int\partial_jT^{ij}(\mathbf{x},t)\,\dd^3x
  =0.
\]
Thus the shift in origin does not affect the Fourier components with
\(\omega\ne0\), that is,
\begin{equation}
  D_{ij}(\omega)
  \equiv
  \int x^ix^jT^{00}(\mathbf{x},\omega)\,\dd^3x
  =
  \int(x^i-a^i)(x^j-a^j)T^{00}(\mathbf{x},\omega)\,\dd^3x.
  \tag{10.5.12}\label{eq:10.5.12}
\end{equation}
However, it is only when we take \(T^{00}\) as the energy density of the
entire system that we can shift origins freely in computing
\(D_{ij}(\omega)\).

As a first example, let us calculate the gravitational radiation produced by
a sound wave in a tube lying in the \(z\)-direction. The density of the
vibrating material can be written
\[
  \rho=\rho_0+\rho_1,
\]
where \(\rho_0\) is the constant unperturbed value and \(\rho_1\) is a
small perturbation. We also treat the material velocity \(v\) (in the
\(z\)-direction) as a small perturbation, and neglect dissipative effects.
The equations of motion are then
\[
  \partial_t\rho_1+\rho_0\partial_zv=0,
  \qquad
  \rho_0\partial_tv+v_s^2\partial_z\rho_1=0,
\]
where \(v_s\) is the speed of sound. The tube is not supported at its ends
(otherwise we would have to take into account the gravitational radiation
emitted by the support!), so the pressure \(v_s^2\rho_1\) must vanish at the
tube ends. With this boundary condition, the general solution for a tube
extending from \(z=0\) to \(z=L\) is a superposition of the normal modes
\begin{equation}
  v=-\epsilon_{\mathrm{osc}} v_s\cos kz\sin(\omega t+\phi),
  \tag{10.5.13}\label{eq:10.5.13}
\end{equation}
\begin{equation}
  \rho_1=\epsilon_{\mathrm{osc}}\rho_0\sin kz\cos(\omega t+\phi),
  \tag{10.5.14}\label{eq:10.5.14}
\end{equation}
where \(\epsilon_{\mathrm{osc}}\) is a small dimensionless oscillation amplitude,
\(\phi\) is an arbitrary
phase, and
\begin{equation}
  k=N\frac{\pi}{L},
  \qquad
  \omega=N\pi\frac{v_s}{L},
  \tag{10.5.15}\label{eq:10.5.15}
\end{equation}
with \(N\) any positive integer. Since \(v\) is not constrained to vanish at
the tube ends, these ends will in general be displaced by amounts
\(\delta(0,t)\) and \(\delta(L,t)\), respectively, where
\[
  \delta(z,t)
  \equiv
  \int v(z,t)\,\dd t
  =
  \epsilon_{\mathrm{osc}}v_s\omega^{-1}\cos kz\cos(\omega t+\phi).
\]
The time-dependent part of the second moment of the mass density is given
here by
\[
  D_{ij}(t)
  =
  n_in_jA
  \left(
    \int_0^L\rho_1(z,t)z^2\,\dd z
    +L^2\rho_0\delta(L,t)
  \right),
\]
where \(A\) is the cross-sectional area of the tube, and
\(\mathbf{n}=(0,0,1)\) is a unit vector in the \(z\)-direction. This
vanishes for \(N\) even, whereas for \(N\) odd we find
\[
  D_{ij}(t)
  =
  -\left(
    \frac{4n_in_jML^2\epsilon_{\mathrm{osc}}}{N^3\pi^3}
  \right)\cos(\omega t+\phi),
\]
where \(M\equiv\rho_0AL\) is the mass of the tube. (The reader can easily
check that \(D_{ij}(t)\) would be the same if the second moment of the mass
distribution were evaluated using as origin some point other than \(z=0\).)
Comparing with Eq.~\eqref{eq:10.5.11}, we see that \(D_{ij}(t)\) has Fourier
component
\begin{equation}
  D_{ij}\!\left(N\pi\frac{v_s}{L}\right)
  =
  -\frac{2n_in_jML^2\epsilon_{\mathrm{osc}}}{N^3\pi^3}.
  \tag{10.5.16}\label{eq:10.5.16}
\end{equation}
The radiated power is thus given by Eq.~\eqref{eq:10.5.7} (in c.g.s.
units) as
\begin{equation}
  P
  =
  \frac{16GM^2v_s^6\epsilon_{\mathrm{osc}}^2}{15L^2c^5},
  \tag{10.5.17}\label{eq:10.5.17}
\end{equation}
for each odd \(N\). This may be compared with the total energy of the
oscillation, which is simply the kinetic energy at the times when \(\rho_1\)
vanishes, that is, when \(v\) is greatest:
\[
  E
  =
  \frac12\rho_0A\int_0^Lv_{\max}^2(z)\,\dd z
  =
  \frac14Mv_s^2\epsilon_{\mathrm{osc}}^2.
\]
Evidently the emission of gravitational radiation will cause the oscillator
to lose energy at a rate
\begin{equation}
  \Gamma_{\mathrm{grav}}
  \equiv
  \frac{P}{E}
  =
  \frac{64GMv_s^4}{15L^2c^5}.
  \tag{10.5.18}\label{eq:10.5.18}
\end{equation}

For instance, let us calculate the rate of gravitational radiation by
acoustic oscillations in the large aluminum cylinders used as antennas in
Weber's experiments on gravitational radiation.%
\hyperref[ch10-ref-1]{\textsuperscript{1}}
(As we shall see, the effective cross-section of such antennas is determined
by \(\Gamma_{\mathrm{grav}}\).) Weber's cylinders have the parameters
\[
  L=153\ \mathrm{cm},
  \qquad
  v_s=5.1\times10^5\ \mathrm{cm/sec},
  \qquad
  M=1.4\times10^6\ \mathrm{gm}.
\]
Hence, if gravitational radiation were the only loss mechanism, the
oscillations \eqref{eq:10.5.13}, \eqref{eq:10.5.14} with \(N\) odd would
lose energy at a rate
\begin{equation}
  \Gamma_{\mathrm{grav}}
  =
  4.7\times10^{-35}\ \mathrm{sec}^{-1}.
  \tag{10.5.19}\label{eq:10.5.19}
\end{equation}
In contrast, the actual decay rate \(\Gamma\) of the \(N=1\) mode in this
cylinder is about \(0.15\ \mathrm{sec}^{-1}\), owing primarily to viscous
dissipation within the aluminum. Hence the ``branching ratio'' of
gravitational radiation here is of order
\begin{equation}
  \eta\equiv\frac{\Gamma_{\mathrm{grav}}}{\Gamma}
  \simeq3\times10^{-34}
  \qquad(N=1).
  \tag{10.5.20}\label{eq:10.5.20}
\end{equation}
Any ordinary mechanical oscillation will always give up vastly more of its
energy to heat than to gravitational radiation.

As a second example, let us calculate the power radiated by a rotating body.
If the body rotates rigidly about the 3-axis with angular frequency
\(\Omega\), then the mass density \(T^{00}\) will take the form
\[
  T^{00}(\mathbf{x},t)=\rho(\mathbf{x}'),
\]
where \(\rho(\mathbf{x}')\) is the mass density expressed in coordinates
\(\mathbf{x}'\) fixed in the body, defined by
\[
  x_1=x'_1\cos\Omega t-x'_2\sin\Omega t,
  \qquad
  x_2=x'_1\sin\Omega t+x'_2\cos\Omega t,
  \qquad
  x_3=x'_3.
\]
Hence, by changing coordinates in \eqref{eq:10.5.9}, we may express
\(D_{ij}(t)\) in terms of the moment-of-inertia tensor in body-fixed
coordinates
\begin{equation}
  I_{ij}
  \equiv
  \int \dd^3x'\,x'_ix'_j\rho(\mathbf{x}').
  \tag{10.5.21}\label{eq:10.5.21}
\end{equation}
For simplicity, let us consider rotation around one of the principal axes of
the ellipsoid of inertia, so that \(I_{13}=I_{23}=0\). We can also choose
the \(x'_1\) and \(x'_2\) axes along the two other principal axes, so that
\(I_{12}=0\). With \(I_{ij}\) diagonal, we now find
\[
  D_{11}(t)
  =\frac12(I_{11}+I_{22})
  +\frac12(I_{11}-I_{22})\cos2\Omega t,
\]
\[
  D_{12}(t)=\frac12(I_{11}-I_{22})\sin2\Omega t,
\]
\[
  D_{22}(t)
  =\frac12(I_{11}+I_{22})
  -\frac12(I_{11}-I_{22})\cos2\Omega t,
\]
\[
  D_{13}(t)=D_{23}(t)=0,
  \qquad
  D_{33}(t)=I_{33}.
\]
The nonvanishing Fourier coefficients for \(\omega=2\Omega\) in
Eq.~\eqref{eq:10.5.11} are then
\[
  D_{11}(2\Omega)
  =
  -D_{22}(2\Omega)
  =
  \ii D_{12}(2\Omega)
  =
  \frac14(I_{11}-I_{22}).
\]
According to Eq.~\eqref{eq:10.5.7}, the total power emitted at twice the
rotation frequency is then (in c.g.s. units)
\begin{equation}
  P(2\Omega)
  =
  \frac{32G\Omega^6I^2e^2}{5c^5},
  \tag{10.5.22}\label{eq:10.5.22}
\end{equation}
where \(I\) and \(e\) are the moment of inertia and equatorial ellipticity,
\begin{equation}
  I\equiv I_{11}+I_{22},
  \tag{10.5.23}\label{eq:10.5.23}
\end{equation}
\begin{equation}
  e\equiv\frac{I_{11}-I_{22}}{I}.
  \tag{10.5.24}\label{eq:10.5.24}
\end{equation}
A body with circular symmetry around the axis of rotation will have \(e=0\),
and therefore will not emit gravitational radiation. (Indeed, this conclusion
does not even depend on the quadrupole approximation, since such a body,
though rotating, has a time-independent energy-momentum tensor.) On the other
hand, for a point mass \(m\) fixed in the rotating coordinate system at a
point \(x'_1=r,\ x'_2=x'_3=0\), the only nonvanishing element of \(I_{ij}\)
will be \(I_{11}=mr^2\), so that \(I=mr^2\) and \(e=1\), and
Eq.~\eqref{eq:10.5.22} gives a radiated power
\begin{equation}
  P(2\Omega)
  =
  \frac{32G\Omega^6m^2r^4}{5c^5}.
  \tag{10.5.25}\label{eq:10.5.25}
\end{equation}
For instance, for the orbital motion of the planet Jupiter, we have
\[
  \Omega=1.68\times10^{-8}\ \mathrm{sec}^{-1},
  \qquad
  m=1.9\times10^{30}\ \mathrm{g},
  \qquad
  r=7.78\times10^{13}\ \mathrm{cm},
\]
and Eq.~\eqref{eq:10.5.25} gives a gravitational radiation power of only
\(5.3\ \mathrm{kW}\), less even than the solar thermal gravitation power
calculated in the last section. At this rate, it would take very much longer
than the age of the solar system to observe any effects of this energy loss
on Jupiter's orbit.

The negligibility of gravitational radiation in celestial mechanics can be
stated in more general terms. For a system consisting of particles with
typical mass \(\overline M\), typical separations \(\bar r\), and typical
velocities \(\bar v\), the power radiated at a frequency \(\omega\) of order
\(\bar v/\bar r\) will be of order [compare Eq.~\eqref{eq:10.5.7}]
\[
  P\sim
  G\left(\frac{\bar v}{\bar r}\right)^6
  \overline M^{\,2}\bar r^4,
\]
or, since \(G\overline M/\bar r\) is of order \(\bar v^2\),
\[
  P\sim\overline M\frac{\bar v^8}{\bar r}.
\]
The typical deceleration \(\bar a_{\mathrm{rad}}\) of particles in the system
owing to this energy loss is given by the power \(P\) divided by the momentum
\(\overline M\bar v\), or
\[
  \bar a_{\mathrm{rad}}\sim\frac{\bar v^7}{\bar r}.
\]
This may be compared with the accelerations computed in Newtonian mechanics,
which are of order \(\bar v^2/\bar r\), and with the post-Newtonian
corrections discussed in the last chapter, which are of order
\(\bar v^4/\bar r\). [Radiation effects enter with an \emph{odd} power of
\(\bar v\) because they represent an irreversible process, as shown by our
use of an outgoing wave solution in Eq.~\eqref{eq:10.4.4}.] Since radiation
reaction is smaller than the post-Newtonian effects by a factor
\(\bar v^3<10^{-12}\), the neglect of radiation reaction in the last section
was perfectly justified. Indeed, if we had strength we could even compute
the post-post-Newtonian accelerations,%
\hyperref[ch10-ref-2]{\textsuperscript{2}}
which are of order \(\bar v^6/\bar r\), without encountering the effects of
gravitational radiation!

The discovery of the pulsars has provided us with a more promising source of
gravitational radiation. As discussed in Section~11.4, pulsars are probably
neutron stars,%
\hyperref[ch10-ref-3]{\textsuperscript{3}}
with masses of the order of one solar mass, radii of the order of
\(10\ \mathrm{km}\), and hence moments of inertia \(I\) of the order of
\(10^{45}\ \mathrm{g\,cm^2}\). A newborn pulsar, formed in a supernova, may
be rotating with \(\Omega\) of order \(10^4\ \mathrm{sec}^{-1}\), so
according to Eq.~\eqref{eq:10.5.22}, it would be emitting gravitational
radiation at a rate of order \(10^{55}e^2\ \mathrm{ergs/sec}\). For
comparison, the total rotational energy of the pulsar would be about
\(10^{53}\ \mathrm{ergs}\), so most of the pulsar's kinetic energy would be
radiated away as gravitational waves%
\hyperref[ch10-ref-4]{\textsuperscript{4}}
within a few years, provided that the equatorial ellipticity \(e\) is greater
than about \(10^{-4}\). This is too large a static ellipticity to be
maintained in the huge gravitational field of a neutron star, but it might be
possible for dynamical effects to produce a mean ellipticity this large,
particularly in the early period before the pulsar has settled down to its
equilibrium configuration. Eventually the pulsar will slow down sufficiently
so that other loss mechanisms, such as magnetic dipole radiation (for which
\(P\propto\Omega^4\)) become more important than gravitational radiation.
